\chapter{Introduction}
\label{sec:intro}

\section{Motivation and Problem Statement}
Digital design is unforgiving: a single wrong bit in a priority encoder, a missing
reset in a state machine, or a width mismatch in a comparison can silently corrupt a
datapath. An engineer writing RTL must also write the testbench that proves the RTL
correct. Verification is widely reported as the largest share of front-end effort,
and it is exactly the part that generic code generation does not solve.

Applying a general-purpose LLM directly to hardware design fails in three concrete
ways:

\begin{enumerate}[label=\textbf{\arabic*.}]
\item \textbf{Toolchain incompatibility.} Models emit SystemVerilog that the free
simulator does not implement. Icarus Verilog 12.0 in \code{-g2012} mode accepts much of
the synthesizable subset (\code{logic}, \code{always_ff}, \code{return}) but rejects
\code{break} and \code{continue}, \code{disable} of a function, the \code{inside}
operator, streaming concatenation, concurrent assertions (\code{property},
\code{sequence}) and class constraints, so code that looks valid to a human fails to
compile.
\item \textbf{No execution awareness.} A language model cannot observe a simulation
cycle, a clock edge, a race, or an assertion. It can only produce plausible text.
\item \textbf{No real diagnosis.} A model that is simply told ``make it pass'' tends
to weaken the testbench, delete checks, or regenerate a different design. That is
not debugging; it is test evasion.
\end{enumerate}

\section{Project Objectives}
This individual EC486 project builds an execution-grounded assistant with the
following objectives:

\begin{itemize}
\item \textbf{Specification first:} parse an arbitrary prompt into a typed
\code{RTLDesignSpec} with explicit ports, widths, reset polarity, and edge cases,
and let the user review and edit it before anything is generated.
\item \textbf{Independent verification:} generate the testbench from the
specification only, never from the RTL body, so the two cannot trivially agree;
optionally let a different model write the testbench.
\item \textbf{Real simulation:} compile and run every candidate with Icarus
Verilog and decide PASS/FAIL from the simulator, not from the model's claim.
\item \textbf{Structured diagnosis:} require \code{symptom}, \code{root_cause},
\code{fix_location}, and \code{confidence}, grounded in the first failing vector.
\item \textbf{Safe repair:} reject fixes that delete checks, apply mechanical
compiler-evidenced repairs deterministically, and stop on no-op or oscillation.
\item \textbf{Measure the evidence:} after a PASS, report testbench strength by
mutation testing, measured toggle coverage and, when a repair rewrote the testbench,
a second-opinion testbench.
\item \textbf{Early physical estimate:} map a passing design to the SkyWater 130\,nm
cell library and report pre-layout timing and area, or refuse with a reason.
\item \textbf{Honest degradation:} never silently replace a user's design with a
curated demo after a live call has succeeded, and never replace a model the user
chose.
\end{itemize}

\section{Contributions}
The final submission contributes: (i) an end-to-end live pipeline validated on an
arbitrary prompt; (ii) a specification-only testbench generator with explicit Icarus
compatibility rules; (iii) a structured diagnosis/fix loop with deterministic
compiler repairs; (iv) an anti-cheat check-count guard and state-hash oscillation
guard; (v) a truthful degraded-mode policy; (vi) a reproducible test suite that
proves the testbenches can actually fail semantic bugs, with a measured mutation score; (vii) a final review that found and fixed a sandbox hole, a server-freeze defect and testbench gaps; (viii) trust features around every PASS: a reviewable specification, an optional separate testbench model, a second-opinion testbench, a per-run strength score and measured toggle coverage; (ix) a pre-layout timing and area estimate with Yosys and OpenSTA on the SkyWater 130\,nm library, checked by gate-level simulation of the mapped netlist; (x) a per-run model menu over three providers (the OpenCode app, Cloudflare Workers AI and NVIDIA NIM); and (xi) a live evaluation on 15 prompts with independent reference testbenches. The design is model-agnostic, so it scales with stronger models (Section~\ref{sec:scale}).

\section{Report Organisation}
Chapter~\ref{sec:related} surveys related work and Chapter~\ref{sec:req} sets out the requirements. Chapter~\ref{sec:arch} presents the architecture and detailed design, Chapter~\ref{sec:method} the methodology, and Chapter~\ref{sec:impl} the implementation and testing strategy. Chapter~\ref{sec:results} reports the measured results. Chapter~\ref{ch:discussion} discusses the engineering findings, the scalability path and the timeline, deployment and cost. Chapter~\ref{sec:conclusion} concludes, states the limitations and lists future work. Appendices give reproduction steps, the complete prompt set and a live-generated testbench.
